\section{总结与延伸}

V5 overview 的价值在于把 Transformer 从一张 architecture diagram 展开为完整生命周期。Embedding/attention 定义计算接口；数据结构和 schedule 决定梯度；CoT 与 search 增加 test-time compute；preference optimization 改变长期行为；agent 将模型放入环境闭环；视觉和 fMRI 展示 token interface 迁移；continual learning 则追问部署后模型如何真正改变。

\subsection{十二条核心结论}

\begin{enumerate}
  \item Transformer block 是现代 AI 系统的基础组件，不是完整产品解释。
  \item Q/K/V 分别承担查询、匹配索引和被聚合内容的角色。
  \item 数据质量、结构、多样性、规模和 schedule 必须联合设计。
  \item 人类学习对比能生成实验假设，但不能直接证明类脑机制。
  \item TinyDialogues 的 mixed result 表明 diverse mixture 常比单一直觉数据更稳健。
  \item Two-phase pretraining 的收益依赖 blend 与 duration，过采样会递减。
  \item CoT、tree、program 和 decomposition 改变 inference computation，而非同一种“推理”。
  \item Preference methods 的核心差异是监督来源、baseline 与 objective assumptions。
  \item Agent self-improvement 需要外部 verification、权限和停止条件。
  \item ViT/VLM/fMRI 的迁移首先依赖正确 token interface 和 objective。
  \item Scaling limits 需要看局部斜率与新瓶颈，不能用口号判断终结或无限延续。
  \item Continual learning 必须同时获得新能力、控制遗忘并说明更新发生在哪里。
\end{enumerate}

\subsection{一张表串起整堂课}

\begin{table}[H]
\centering
\small
\begin{tabular}{@{}L{0.17\textwidth}L{0.20\textwidth}L{0.28\textwidth}L{0.25\textwidth}@{}}
\toprule
层次 & 主要对象 & 典型方法 & 验收问题 \\
\midrule
Architecture & token interaction & attention, ViT, cross-attention & 表示和复杂度是否合理 \\
Data & 训练分布 & TinyDialogues, two-phase blend & 多指标、重复率、scale 是否稳定 \\
Inference & test-time compute & CoT, ToT, tools & 预算、faithfulness、验证 \\
Preference & 参数行为 & RLHF, DPO, GRPO, KTO & judge、reward、群体差异 \\
Agent & 环境闭环 & ReAct, reflection, LATS & 权限、outcome、恢复 \\
Adaptation & 部署后更新 & editing, MoE, continual FT & plasticity、forgetting、rollback \\
\bottomrule
\end{tabular}
\caption{Transformer 系统从计算到持续适应的六层地图。}
\end{table}

\subsection{实践研究清单}

\begin{importantbox}{开始一个新 Transformer 项目前先回答}
\begin{enumerate}
  \item 输入怎样 tokenized，哪些关系需要 attention？
  \item 目标能力在训练数据中由什么信号监督？
  \item 改进来自更多参数、更多 token、不同 schedule 还是 test-time search？
  \item 是否使用 human/AI feedback，谁来审计 judge？
  \item 工具、memory 与 environment 的权限边界是什么？
  \item 评测是否覆盖 distribution shift、成本、calibration 和失败恢复？
  \item 部署后的更新是外部 memory、局部 editing 还是权重学习？
  \item 新能力是否以 catastrophic forgetting 为代价？
\end{enumerate}
\end{importantbox}

\subsection{复现记录的最小集合}

无论复现 data、reasoning、preference、agent 或 continual-learning 方法，都应保存 source/version、数据 hash、随机种子、模型与 tokenizer、optimizer/schedule、总训练和推理 compute、完整评测 protocol、失败样本以及人工判断规则。若使用 AI judge、工具或检索库，还要固定其版本和权限。只有这些信息齐全，图中的提升才能被另一团队重放、归因和审计；否则 overview 中再漂亮的曲线也只能作为线索，不能成为工程决策依据。

\subsection{自测问题}

\begin{enumerate}
  \item Q、K、V 的矩阵形状和语义分别是什么？为什么要除以 $\sqrt{d_k}$？
  \item Static embedding 与 contextual representation 的差别在哪里？
  \item TinyDialogues 为什么没有支持“纯儿童语料最好”的简单结论？
  \item Two-phase pretraining 中 blend 与 duration 为什么必须分开消融？
  \item CoT、ToT、PoT 与 computation graph 各自增加什么资源？
  \item RLHF、DPO、RLAIF、GRPO、KTO 的监督与基础设施差异是什么？
  \item Self-refinement 在没有外部 verifier 时为什么可能强化错误？
  \item ViT、CLIP 和 VLM 为什么不是同义词？
  \item fMRI foundation model 的 masked objective 学到什么，又不能证明什么？
  \item Scaling 局部斜率下降与“scaling law 失效”有什么区别？
  \item RAG、model editing 和 continual fine-tuning 分别更新什么状态？
  \item 如何同时度量 continual learning 的 plasticity 与 forgetting？
\end{enumerate}

\subsection{拓展阅读}

\begin{itemize}
  \item Vaswani et al., \emph{Attention Is All You Need}：Q/K/V、multi-head attention 与 encoder-decoder。
  \item Dosovitskiy et al., \emph{An Image is Worth 16x16 Words}：Vision Transformer。
  \item Wei et al., \emph{Chain-of-Thought Prompting Elicits Reasoning in Large Language Models}。
  \item Rafailov et al., \emph{Direct Preference Optimization}：preference pairs 与 policy objective。
  \item Yao et al., \emph{ReAct} 与 \emph{Tree of Thoughts}：reasoning/action 和 search。
  \item Meng et al., \emph{ROME}；Mitchell et al., \emph{MEMIT}：model editing。
  \item Lecture slides 中的 TinyDialogues 与 Two-Phase Pretraining primary papers：数据结构与 schedule 的两类实验。
\end{itemize}

\end{document}
